\section{The primitive theta coefficient and integral degeneration scalars}
\label{sec:faithful-psi-normalisation}
We now print the coefficient argument in
\srcid{ExistsDegDataPsiOfSectionNonzeroFibres}, with the same
preuniformisation, complete-discrete-valuation-ring comparison and
nonzero theta section. The argument proves integrality, strict
positivity and adic finiteness of the actual degeneration scalars.
It retains the earlier expansion-existence and generator-comparison
producers specified below; it does not infer their existence from a
formal Fourier series.

Let \(R\) be a normal Noetherian local domain, complete for the
source's proper prime ideal \(I\), and let \(K=\operatorname{Frac}(R)\).
Let \(R\hookrightarrow V\) be the actual injective comparison with a
complete discrete valuation ring, with
\[
                        \mathfrak m_V\cap R=I.
 \tag{PN1}\label{eq:psi-normalisation-centre}
\]
Retain the actual preliminary data \(D,U\) over \(R\), actual full
degeneration datum \(D'\) and uniformisation over \(V\), and the
comparison \(C\) of their abelian quotients, cubical lines, character
lattices, coordinates and rigidifications.
The character-lattice map of \(C\) is the specified isomorphism
\(X'\simeq X\).

\subsection{The actual expansion and generator comparisons}
Take the section \(\sigma\) supplied by
\(\operatorname{ThetaNZ}\), nonzero on every fibre of the actual
family. Its complete-base expansion is \(f=(f_\alpha)_{\alpha\in X}\)
at line power one.
Use the actual comparison to obtain its expansion \(f'\) over
\(V\); the two expansions satisfy the source's
\(\operatorname{SecCmp}\) with the same quotient map, generating
line comparison and lattice isomorphism.
The full uniformisation over \(V\) gives the actual relation
\(\operatorname{PsiRel}(D',f')\).

There is a nonzero term of \(f'\): the finite expansion at the
closed fibre is a unit at some local point, and a finite sum of
nonunits in a local ring cannot be a unit.
The \(\operatorname{PsiRel}\) translations then make every
term \(f'_\alpha\) nonzero. Indeed the translations and their
specified line factors are invertible, and the character-lattice
translation relating any two labels is a bijection.
Every term \(f_\alpha\) is consequently nonzero too.
If all local coefficients of one of them were zero, the
\(\operatorname{SecCmp}\) equality would make every local
coefficient of the corresponding \(f'_\alpha\) zero: its other
factor is a generating line coefficient, hence a unit.
This is \srcid{ne_zero_of_secCmp}, and uses the actual
section comparison rather than just a nonzero dimension count.

The source's
\srcid{exists_psi_psiRel} now supplies the new generic generators
\(\psi_\chi\) with \(\operatorname{PsiRel}(D,f)\).
This generator-construction result is retained here as a separate
input; its conclusion includes the relation on every test scheme.
The rigidification comparison and the same nonzero zero-label
term identify the pullbacks of these generators with those of
\(D'\). Thus one has the actual \(\operatorname{PsiCmp}\).
Over the fraction fields the comparison is along an embedding:
the map \(R\to V\) is injective, so its extension
\(K\to\operatorname{Frac}(V)\) exists uniquely.
The symmetry and bimultiplicative identities for the pulled-back
generators therefore descend by injectivity.
In particular the new preliminary datum has the actual
\(\operatorname{PsiEven}\) and \(\operatorname{TauBimult}\).
The source's line and lattice comparisons then give
\(\operatorname{IsPullDeg}\) with exactly the original \(C.\beta\),
\(C.v\) and \(C.xe\); these equalities remain part of the
conclusion, not just an unspecified isomorphism.

\subsection{The scalar coordinates and the exact coefficient ratio}
For every \(\chi\), choose an integral generating section
\[
                 g_\chi\quad\hbox{of}\quad
                      (c_\chi^*M)^{-1}.
\]
This is possible over the local ring \(R\), since an invertible
module over a local ring is free of rank one.
Write the actual generic generator as
\[
                          \psi_\chi=a_\chi g_\chi,
                          \qquad a_\chi\in K^\times.
 \tag{PN2}\label{eq:psi-normalisation-scalar}
\]
Ideal membership for the generator is then exactly
\[
 \operatorname{PsiIn}(\chi,J)
   \ \Longleftrightarrow\
     a_\chi\in\operatorname{image}(J\longrightarrow K)
               \qquad(J\subset R).
 \tag{PN3}\label{eq:psi-normalisation-ideal}
\]
To verify this equivalence in every chart, the coefficient of
\(g_\chi\) is a unit, so membership of its multiple in the
extended ideal is equivalent to membership of the scalar.
Conversely evaluate the global rank-one trivialisation over
\(\operatorname{Spec}R\) to recover the scalar.
Changing an integral generator multiplies the scalar by an
integral unit and preserves the ideal.
This is the exact local content of
\srcid{PsiScalIn.main_local}.

Let \(\vartheta\) be the source's nonzero-fibre theta section on
the abelian quotient. Its line-power-one comparison supplies
\(q,b_\chi\in R\) with
\[
                  f_0=q\vartheta,\qquad
                  \operatorname{Coef}
                     (D,\chi,g_\chi,\vartheta,f_\chi,b_\chi).
 \tag{PN4}\label{eq:psi-normalisation-coefficients}
\]
The second relation is the actual translated-section coefficient
identity; its line factors are the prescribed ones.
It is not a choice of an arbitrary coefficient functional.

\begin{srlemma}[The source coefficient ratio]
\label{lem:psi-normalisation-ratio}
For all \(\chi\),
\[
                             q a_\chi=b_\chi
                                      \quad\hbox{in }K.
 \tag{PN5}\label{eq:psi-normalisation-ratio}
\]
\end{srlemma}
\begin{proof}
Choose a point in the nonzero locus of the generic fibre of
\(\vartheta\), and work in its local ring.  Its theta coefficient
\(t\) is a unit.  Evaluate the
\(\operatorname{PsiRel}\) identity after translating by
\(-c_\chi\), so that the subsequent translation by \(c_\chi\)
returns to that same point. In consistent frames let
\(s,s_0\) denote the coefficients of the two expansion terms,
\(u\) the actual \(\psi_\chi\) coefficient, \(g\) its integral
generator coefficient, and \(r\) the remaining prescribed
invertible cubical factor.
The four defining identities are precisely
\[
             s=ur\,s_0,\qquad u=a_\chi g,\qquad
             s_0=qt,\qquad s=b_\chi(gr)t .
\]
All scalars are mapped from \(K\) into this same local ring.
Substitution and cancellation of the unit \(grt\) give
\(qa_\chi=b_\chi\) there. The map from the field \(K\) into a
nonzero local \(K\)-algebra is injective, so the equality holds
in \(K\).  This expands
\srcid{psiScal_eq_ratio}, including the common stalk and the
specific translating point needed for its cancellation.
\end{proof}

\begin{srlemma}[Positivity transported before division by \(q\)]
\label{lem:psi-normalisation-positive-coefficients}
For every nonzero \(\chi\in X\), \(b_\chi\in I\).
\end{srlemma}
\begin{proof}
Let \(\chi'\) be its preimage under \(C.xe\). The full datum over
\(V\) has \(\operatorname{PsiIn}(\chi',\mathfrak m_V)\).
Choose its integral rank-one generator as in
\eqref{eq:psi-normalisation-scalar}.
Pulling \(g_\chi\) through the actual generating line
comparison gives another integral generator on \(V\).
The ratio of these two integral generators is a unit of \(V\).
The equality \(\operatorname{PsiCmp}\) therefore says that the
image of \(a_\chi\) in \(\operatorname{Frac}(V)\) is
\(m_\chi\in\mathfrak m_V\): it is the product of the
coefficient in \(\mathfrak m_V\) and that integral transition
unit.
Apply the fraction-field embedding to
\eqref{eq:psi-normalisation-ratio}. It gives the integral
identity
\[
                         \phi(b_\chi)=\phi(q)m_\chi .
\]
The right side lies in \(\mathfrak m_V\), since \(q\in R\).
Equation~\eqref{eq:psi-normalisation-centre} gives
\(b_\chi\in I\).
This is the producer
\srcid{scalar_mem_of_psiIn}. It requires no preliminary
assumption that \(q\) is invertible.
\end{proof}

\subsection{Primitivity forces the zero-label coefficient to be a unit}
\begin{srlemma}[An actual unit term on the closed fibre]
\label{lem:psi-normalisation-unit-term}
There are an actual test scheme \(T/R\), a point \(z\in T\)
above the closed point of \(R\), compatible quotient frames,
and a label \(\alpha\), such that the evaluated coefficient
of \(f_\alpha\) is a unit at \(z\).
Every element of \(I\) maps to zero at that same test point.
\end{srlemma}
\begin{proof}
The section \(\sigma\) is nonzero on the closed fibre.
Choose a point in its nonzero locus and pull it to the
level-zero thickening \(G_0=G\times_R R/I\).
Such a pullback point exists since \(I\) is contained in
the maximal ideal of the local base.
Use the actual torus quotient \(q_0\) and the actual
generating line comparison of \(U\).
The level-zero expansion is finite. In its common local
frames it writes the unit value of \(\sigma\), multiplied
by a generating unit, as a finite sum of coefficients of
the \(f_\alpha\), multiplied by character units.
If every summand were a nonunit, all would lie in the
maximal ideal of this local ring, and so would their sum.
Thus a summand is a unit, and hence its \(f_\alpha\)
coefficient is a unit.
Take \(T=G_0\) with its structural map to \(\operatorname{Spec}R\).
It factors through \(R/I\), so all elements of \(I\) map
to zero.  All these assertions use that same \(T,z\) and
the same actual quotient frames.
This is \srcid{exists_unit_term}.
\end{proof}

\begin{srproposition}[The primitive denominator is a unit]
\label{prop:psi-normalisation-q-unit}
The scalar \(q\) in \eqref{eq:psi-normalisation-coefficients}
belongs to \(R^\times\). Consequently every \(a_\chi\)
belongs to \(R\), and \(a_\chi\in I\) for \(\chi\ne0\).
\end{srproposition}
\begin{proof}
Use Lemma~\ref{lem:psi-normalisation-unit-term}.
If its label \(\alpha\) is nonzero,
Lemma~\ref{lem:psi-normalisation-positive-coefficients}
makes \(b_\alpha\in I\).
The actual \(\operatorname{Coef}\) identity at the same point
then evaluates \(f_\alpha\) to zero, since its scalar
\(b_\alpha\) is zero there. This contradicts its being a unit.

Hence \(\alpha=0\). Equation \(f_0=q\vartheta\) makes the image
of \(q\) a unit in this local ring, since a factor of a
unit product is a unit.
Its structural map lies above the closed point of
\(\operatorname{Spec}R\), so a nonunit of \(R\), belonging
to \(\mathfrak m_R\), would map into the local maximal
ideal.  Therefore \(q\) was already a unit in \(R\).
Now \eqref{eq:psi-normalisation-ratio} gives
\(a_\chi=q^{-1}b_\chi\in R\), and its positive-label
ideal assertion follows because multiplication by a unit
preserves \(I\).
\end{proof}

\subsection{The uniform adic finiteness assertion}
\begin{srlemma}[The closed-fibre flat-stalk test]
\label{lem:psi-normalisation-flat-stalk}
If \(S\) is local, \(p:Z\to\operatorname{Spec}S\) is flat
and \(z\) lies above its closed point, then
\(S\to\mathcal O_{Z,z}\) is injective.
\end{srlemma}
\begin{proof}
It is a flat local homomorphism and hence faithfully flat.
A faithfully flat algebra has injective structural map:
if \(s\) maps to zero, the inclusion of the cyclic ideal
\((s)\) in \(S\) becomes zero after faithfully flat
extension, and faithfulness forces \((s)=0\).
This is \srcid{eq_zero_of_cst_eq_zero}.
\end{proof}

\begin{srproposition}[The actual expansion coefficients tend adically to zero]
\label{prop:psi-normalisation-adic}
For every \(e\geq0\) there is a finite set \(S_e\subset X\)
such that \(b_\chi\in I^{e+1}\) whenever \(\chi\notin S_e\).
\end{srproposition}
\begin{proof}
At level \(e\), the actual expansion condition gives a
finite set \(S_e\) outside which the \(f_\chi\) coefficients
vanish on \(G_e=G\times_R R/I^{e+1}\).
Fix \(\chi\notin S_e\).
Choose a point of the closed fibre of the abelian quotient
at which \(\vartheta\) is nonzero.
Translate it by \(-c_\chi\), and lift it to \(G_e\)
through the actual surjective torus quotient.
In the notation of the coefficient identity, the
subsequent \(+c_\chi\) translation returns to the chosen
nonzero-theta point. Thus at the lifted point \(z\) the
coefficient of \(f_\chi\) is zero, while the product of
the integral \(g_\chi\) coefficient, the prescribed
cubical factors and the translated theta coefficient
is a unit.
The \(\operatorname{Coef}\) identity therefore makes the
image of \(b_\chi\) zero in \(\mathcal O_{G_e,z}\).

The base ring \(R/I^{e+1}\) is local, and \(z\) lies
above its closed point. The map from \(G_e\) is smooth,
hence flat, since it is a pullback of the actual
semiabelian family. Lemma~\ref{lem:psi-normalisation-flat-stalk}
therefore detects the coefficient:
\(b_\chi\bmod I^{e+1}=0\).
This proves the stated ideal membership.
It is the exact same-point argument of
\srcid{coef_mem_pow}; ordinary generic-point vanishing
would not detect this integral ideal.
\end{proof}

\begin{srcorollary}[The positivity and finiteness fields of the actual datum]
\label{cor:psi-normalisation-degdata}
For the new \(\psi_\chi\),
\[
 \operatorname{PsiIn}(\chi,I)\quad(\chi\ne0),\qquad
 \{\chi:\neg\operatorname{PsiIn}(\chi,I^n)\}\ \hbox{is finite}
                                      \quad(n\geq0).
\]
Together with the descended \(\operatorname{PsiEven}\)
and \(\operatorname{TauBimult}\), these are the added
fields of the actual \(\operatorname{DegData}(R,I)\).
The underlying preliminary objects and their
comparison with \(D'\) are retained.
\end{srcorollary}
\begin{proof}
Use \(a_\chi=q^{-1}b_\chi\) and
\eqref{eq:psi-normalisation-ideal}.
For finiteness, use Proposition~\ref{prop:psi-normalisation-adic}
with \(e=n\) and
\(I^{n+1}\subset I^n\); the exceptional labels lie in the
finite set \(S_n\).  For \(n=0\), this argument also works
with \(I^0=R\).
Replacing \(\psi\) changes no field in the source's
\(\operatorname{PreU}\) that records the quotient maps,
coordinates, generic polarized variety or line on \(G\).
The source transports those fields literally by
\srcid{preWithPsi} and its comparison by
\srcid{cmpWithPsi}. Thus the constructed datum has the
specified original comparison maps.
\end{proof}

\paragraph{Precise source boundary.}
The four elementary coefficient and ideal producers
proved here are, at the frozen commit
\srcid{7cbf7a0c32fca1badf05931879a37c3f8a2b03b4},
\srcid{Development/AbelianVarietyTheory/Proofs/DegenerationProofs/PsiScalarEqExpansionCoefRatio.lean},
\srcid{PsiScalarMemOfPsiMemIdeal.lean},
\srcid{ExistsExpansionTermUnitClosedFibre.lean}, and
\srcid{ExpansionCoefMemAdicPow.lean}.
The opening assembly is
\srcid{ExistsDegDataPsiOfSectionNonzeroFibres.lean}.
It retains the actual complete-base expansion-existence
theorem, fibrewise nonzero-section base change,
the construction \srcid{exists_psi_psiRel}, its rigidification
and generator comparisons, and their prescribed
\(\operatorname{IsPullDeg}\) equalities.
This chapter does not replace those producers by the
primitive coefficient argument. The separate descent
of the full expansion bijections remains necessary
for the complete-local uniformisation in N3.
